\subsection{RQ3: What interaction Techniques enable high-fidelity XR prototyping?}

This study employs the Mixed Reality Toolkit (MRTK) in combination with the Unity XR Toolkit to scaffold the user interface. Control inputs are implemented using the Meta Quest Touch Controller to emulate real-world quadcopter interaction. The user operates within a local-floor reference space, while a view-based heads-up display (HUD) provides real-time control status feedback.

To simulate first-person-view (FPV) goggles, Unity's RenderTexture is used to project the quadcopter's camera feed onto a virtual display panel. This design enables users to simultaneously perceive both the surrounding environment and the drone's onboard perspective, thereby supporting situational awareness and potentially reducing simulator sickness.

The system incorporates three interaction environments: (1) Additive Environment Blend, an augmented reality setting in which users perform tutorial tasks within their physical surroundings; (2) a white-boxed playground, which serves as a controlled training environment for developing piloting skills; and (3) real city piloting, a virtual geographic environment that enables navigation within a simulated urban landscape. The design of these learning environments is grounded in experiential learning theory \parencite{kolb1984}, which conceptualizes learning as a cyclical process of experience, reflection, and experimentation.

\paragraph{When XR helps}

This study positions XR as a spatial prototyping medium for FPV drone learning system design. XR enables direct manipulation of control mappings in a three-dimensional environment, reducing the need for abstract interface layout design and accelerating iterative testing of interaction strategies.

\paragraph{When XR is not preferable}

When spatial arrangement is not part of the design problem, XR may introduce unnecessary complexity without improving simulation fidelity or development efficiency. In such cases, traditional simulation or parameter optimization approaches, such as PID tuning or reinforcement learning-based control refinement, are more efficient.
